% !TEX root = CoeneDylanBP.tex
%%=============================================================================
%% Samenvatting
%%=============================================================================

% TODO: De "abstract" of samenvatting is een kernachtige (~ 1 blz. voor een
% thesis) synthese van het document.
%
% Een goede abstract biedt een kernachtig antwoord op volgende vragen:
%
% 1. Waarover gaat de bachelorproef?
% 2. Waarom heb je er over geschreven?
% 3. Hoe heb je het onderzoek uitgevoerd?
% 4. Wat waren de resultaten? Wat blijkt uit je onderzoek?
% 5. Wat betekenen je resultaten? Wat is de relevantie voor het werkveld?
%
% Daarom bestaat een abstract uit volgende componenten:
%
% - inleiding + kaderen thema
% - probleemstelling
% - (centrale) onderzoeksvraag
% - onderzoeksdoelstelling
% - methodologie
% - resultaten (beperk tot de belangrijkste, relevant voor de onderzoeksvraag)
% - conclusies, aanbevelingen, beperkingen
%
% LET OP! Een samenvatting is GEEN voorwoord!

%%---------- Nederlandse samenvatting -----------------------------------------
%
% TODO: Als je je bachelorproef in het Engels schrijft, moet je eerst een
% Nederlandse samenvatting invoegen. Haal daarvoor onderstaande code uit
% commentaar.
% Wie zijn bachelorproef in het Nederlands schrijft, kan dit negeren, de inhoud
% wordt niet in het document ingevoegd.

\IfLanguageName{english}{%
    \selectlanguage{dutch}
    \chapter*{Samenvatting}
    \lipsum[1-4]
    \selectlanguage{english}
}{}

\chapter*{\IfLanguageName{dutch}{Samenvatting}{Abstract}}

Trampolinespringen is een technische turndiscipline waarbij atleten worden beoordeeld op onder meer horizontale verplaatsing (HD) en vluchttijd (ToF). In officiële competities worden deze waarden objectief gemeten met het HDTS Trampoline Measurement Device, een systeem gebaseerd op krachtplaten dat tussen \euro{}5.000 en \euro{}10.000 kost. Die hoge kostprijs maakt het systeem ontoegankelijk voor veel turnclubs, die daardoor aangewezen zijn op tijdsintensieve, manuele videoanalyse.

Dit onderzoek stelt de vraag of een goedkoop computer vision-systeem een betrouwbaar alternatief kan bieden: \textit{Hoe kan een goedkoop computer vision-systeem worden ontwikkeld dat de scores voor horizontale verplaatsing en vluchttijd bij trampolinespringen betrouwbaar kan benaderen, als betaalbaar alternatief voor het HDTS-meetsysteem?} Het doel is een werkend proof-of-concept dat aantoont dat HD-landingszones automatisch kunnen worden onderscheiden en vluchttijden kunnen worden geschat op basis van videobeelden alleen.

De ontwikkelde pipeline is iteratief opgebouwd en bestaat uit vier componenten: hoekpuntdetectie via een gefinetuned YOLO11n-pose model, een homografie die pixelposities omzet naar metrische waarden op het trampolinebed, een binaire classifier voor het detecteren van het landingsmoment, en SAM-2 segmentatie voor de exacte landingspositie. Alternatieve benaderingen, waaronder pose estimation met ViTPose en optische stroming, werden experimenteel geëvalueerd en gemotiveerd verworpen. De pipeline werd geëvalueerd op 49 testvideo's uit de publieke dataset van \textcite{Lin2025}.

De resultaten tonen dat de pipeline HD-landingszones betrouwbaar onderscheidt in 1D-modus, met een gemiddelde absolute afwijking (MAE) van 0,649 ten opzichte van de cumulatieve HDTS-score over een volledige reeks. De ToF-schatting via frame-telling levert een MAE van 5,253 seconden op, wat de gestelde eis van 0,2 seconden niet haalt. De voornaamste beperkingen zijn de lage framerate van de gebruikte wedstrijdvideo's en de gevoeligheid van de homografietransformatie voor fouten langs de diepte-as.

De conclusie is dat een low-cost computer vision-pipeline de fundamentele componenten van HD- en ToF-meting technisch kan realiseren. Het systeem draait volledig op open-source modellen en lichtgewicht hardware, wat de potentiële kostprijs significant lager maakt dan die van het HDTS-systeem. De huidige nauwkeurigheid is nog onvoldoende voor directe inzet in competitieverband, maar de resultaten tonen een geloofwaardig startpunt voor verder onderzoek. Met een hogere framerate, een tweede camera voor de diepte-as en een dataset met per-sprong HDTS-annotaties zou een significant nauwkeuriger systeem realiseerbaar zijn.
